% Generated draft: requires mathematical and source review.
\documentclass[11pt]{article}
\usepackage[margin=25mm]{geometry}
\usepackage{fontspec}
\setmainfont{Latin Modern Roman}
\usepackage{amsmath,amssymb,mathtools,unicode-math}
\setmathfont{Latin Modern Math}
\usepackage{enumitem,needspace,xurl,hyperref}
\setlength{\parindent}{0pt}
\setlength{\parskip}{5pt}
\setlength{\emergencystretch}{3em}
\newcommand{\N}{\mathbb{N}}
\newcommand{\Z}{\mathbb{Z}}
\newcommand{\Q}{\mathbb{Q}}
\newcommand{\R}{\mathbb{R}}
\newcommand{\C}{\mathbb{C}}
\newcommand{\F}{\mathbb{F}}
\newcommand{\A}{\mathbb{A}}
\newcommand{\PP}{\mathbb{P}}
\newcommand{\E}{\mathbb{E}}
\newcommand{\eps}{\varepsilon}
\newcommand{\dd}{\,\mathrm{d}}
\begin{document}
\section*{UnsolvedMath ID: 30001580}
\textbf{Source problem code:} catalogue entry
\subsection{Definitions and mathematical statement}
For a finitely generated group \(G\), define its rank by
\[
d(G)=\min\{\lvert S\rvert:S\subseteq G,\ \langle S\rangle=G\}.
\]
A closed hyperbolic \(3\)-manifold is a compact \(3\)-manifold without boundary equipped with a complete Riemannian metric of constant sectional curvature \(-1\). Connectedness is implicit in the source's use of a single fundamental group.

For a connected, closed, orientable \(3\)-manifold \(M\), a Heegaard splitting of genus \(h\) is a decomposition
\[
M=H_1\cup_{\Sigma}H_2,
\]
where each \(H_i\) is a handlebody homeomorphic to a regular neighborhood of a bouquet of \(h\) circles, and \(H_1\cap H_2=\partial H_1=\partial H_2=\Sigma\). Its Heegaard genus is
\[
g(M)=\min\{h\in\Z_{\geq0}:M\text{ admits a Heegaard splitting of genus }h\}.
\]

The supplied question asks whether
\[
d\bigl(\pi_1(M)\bigr)=g(M)
\qquad\text{for every closed hyperbolic }3\text{-manifold }M.
\]
The source does not explicitly impose orientability. The definition of Heegaard genus above gives the usual orientable formulation; whether the source intends only orientable manifolds or a convention extending Heegaard genus to nonorientable manifolds remains unresolved. Thus the full domain of the supplied universal question requires clarification.
\subsection{Short English statement}
Must the minimum number of generators of the fundamental group of every closed hyperbolic \(3\)-manifold equal its Heegaard genus? The source leaves orientability implicit, so the intended manifold class and any nonorientable genus convention require review.
\subsection{Sources}
\begin{itemize}
\item Local UnsolvedMath registry, permanent ID \(30001580\), source problem code \texttt{catalogue entry}, title ``Rank Versus Heegaard Genus of Hyperbolic Three-Manifolds.'' The supplied proposed year is 2010; no individual proposer is identified.
\item Supplied source citation: \emph{Actions and Invariants of Residually Finite Groups: Asymptotic Methods} (2011), \emph{Oberwolfach Reports}, pp.\ 2335--2391, DOI 10.4171/owr/2010/40. Original OWR report: \url{https://doi.org/10.4171/owr/2010/40}.
\item Supplied literature source: Aschenbrenner--Friedl--Wilton, \emph{3-manifold groups}, \url{https://www.maths.gla.ac.uk/~mpowell/Aschenbrenner-Friedl-Wilton-3-manifold-groups-final-version-031115.pdf}.
\item Literature assessment dated \texttt{2026-08-21}: the supplied curation labels this question open, states that the general statement fails, and says that the report isolates the hyperbolic three-manifold case as open and that later surveys continue to list corresponding rank and mod-\(p\) homology questions as unresolved. This is supplied dated curation, not independently verified here; the broader statement and the mod-\(p\) questions are not specified in the record.
\item The registry marks the statement \texttt{corrected\_verified} and reports comparison with the official Oberwolfach report, with notation, context, and question form normalized without changing the source-level claim. These are supplied verification assertions, not independent verification in this draft.
\end{itemize}
\end{document}
